\section{W5, W6, and the Gray-code special case}\label{sec:special}

The two nontrivial affine operators describe very different state semantics.  Along one context
class with initial state $q$, Equations~\eqref{eq:w5-rule} and \eqref{eq:w6-rule} unwind to
\begin{align}
\W5:\quad
  r_0&=y_0+q,
  &r_j&=y_j+y_{j-1} &&(j\ge1),\label{eq:w5-unwind}\\
\W6:\quad
  r_j&=q+\sum_{i=0}^{j}y_i .&&&&\label{eq:w6-unwind}
\end{align}
$\W5$ asks whether the current successor differs from the last successor of this context.  Its
state is a prediction with a clear learning rule.  $\W6$ reports cumulative parity; because its
new state equals the emitted $r_j$, it feeds the residual back into the next visit.  Its state
continues to contain the initial bit and all previous successors.

\begin{figure}[ht]
  \centering
  \begin{tikzpicture}[
    >={Stealth[length=2.2mm]},
    box/.style={draw, rounded corners=1.5pt, inner xsep=6pt, inner ysep=4pt, font=\small},
    xorg/.style={draw, circle, minimum size=5mm, inner sep=0pt},
    lbl/.style={font=\footnotesize},
    wire/.style={line width=0.5pt},
    fb/.style={line width=1.1pt},
  ]
  \node[box, minimum width=1.6cm] (S) at (0,0) {$\cell{\addr{k}}$};
  \node[xorg] (xr) at (2.2,0) {$+$};
  \node[box] (out) at (3.9,0) {$\out(k)$};
  \draw[->,wire] (S.east) -- node[above,lbl] {$u$} (xr.west);
  \draw[->,wire] (2.2,1.25) -- node[right,lbl,pos=0.15] {$\tape(k)$} (xr.north);
  \draw[->,wire] (xr.east) -- (out.west);
  \fill (2.2,0.8) circle (1.5pt);
  \draw[->,fb] (2.2,0.8) -- (-1.45,0.8) -- (-1.45,-1.0) -- (0,-1.0) -- (S.south);
  \node[lbl] at (1.2,-1.5) {$\W5:\; \cell{\addr{k}}\leftarrow\tape(k)$};
\end{tikzpicture}
\hspace{1.1cm}\input{fig/wiring6.tex}
  \caption{$\W5$ and $\W6$ use the same emitter.  Moving the feedback tap from before the
  exclusive-or to after it changes a last-observation predictor into a running-parity state.}
  \label{fig:w5-w6-taps}
\end{figure}

The operators $1+T$ and $(1+T)^{-1}$ are mutual inverses, but this statement has a precise scope.
If the address sequence is held fixed, applying the $\W6$ recurrence to the $\W5$ residual (or
vice versa) restores each per-context sequence.  In $\cxor$, however, addresses are generated
from the transform's current input.  Feeding a $\W5$ residual into a fresh $\W6$
encoder generally produces a different address sequence and is not the decoder of
Section~\ref{sec:cta}.  The causal decoder reconstructs source bits and reproduces the original
addresses while applying the \emph{same} update rule.  Operator duality and transform inversion
are therefore related but distinct facts.

\subsection{Order zero}

At $m=0$ there is only one context, so visits are ordinary adjacent positions and $T$ is the
usual one-step delay.  With a zero initial cell, $\W5$ gives
\begin{equation}\label{eq:gray}
  r_0=x_0,
  \qquad
  r_k=x_k+x_{k-1}\quad(k\ge1).
\end{equation}
If the bits are read most significant first as the binary representation of an integer $b$,
Equation~\eqref{eq:gray} is exactly reflected binary (Gray) encoding
\[
  g=b\mathbin{\oplus}(b\mathbin{\gg}1)
\]
\cite{gray1953}.  $\W6$ takes prefix exclusive-ors and is the corresponding Gray decoder.  This
is the simplest concrete instance of the two transfer operators.

For example, $90$ is $\bstr{01011010}$ in eight-bit binary.  $\W5$ emits
$\bstr{01110111}$, representing $119=90\oplus45$, and $\W6$ maps that word back to
$\bstr{01011010}$.  In this order-zero case the fixed-address operator duality is also an ordinary
composition of the two transforms, because there is no source-dependent address sequence to
change between passes.

For $m>0$, the delay means ``previous visit to this context'' rather than ``previous tape
position.''  At $m=1$ with zero table and prefix, $\W5$ happens to reduce to the stride-two
difference $r_k=x_k+x_{k-2}$, taking out-of-range bits as zero.  At larger orders the visit
partition is source-dependent and no fixed tape stride describes it.  Calling $\cxor$ a
generalised Gray code would therefore hide its defining feature: context-addressed prediction.
Gray coding is an elegant boundary case, not the main interpretation.
